\chapter{经典CSMA饱和吞吐量的统一参数推导}
\label{app:classic_csma}

本附录基于休假排队分析框架，重新推导经典CSMA的饱和吞吐量。
推导采用与正文一致的参数体系。
结果可作为与 CSMAST~比较的理论基准。

\section{统一参数表示与模型设定}

在已有文献中，例如 \cite{2013_DaiCSMA}，经典CSMA的最大吞吐量常以传播时延参数 $a$ 表示。
该文献采用状态时长设定
$$
\tau_S=1+a,\quad \tau_F=(x+1)a,\quad \tau_I=a.
$$
由此得到的最大吞吐量表达式包含 Lambert W 函数
\begin{equation}
	\hat{\lambda}_{\text{max}} = \frac{-W_{0}\left(-\frac{1}{x(1+1/x)}\right)}{xa-(1-xa)\,W_{0}\left(-\frac{1}{x(1+1/x)}\right)}.
\end{equation}
为统一比较维度，本文不固定参数 $a$。
我们直接使用状态时长参数 $\tau_S$、$\tau_F$ 与 $\tau_I$ 进行推导。

在本章模型下，经典 CSMA 可视为 CSMAST 的特例。
两者关键差异在于忙期构造方式。
在 CSMAST 中，节点若未发生碰撞即可连续发送，因此忙期内成功发送次数服从参数为 $1-p_0$ 的几何分布，其均值为
$$
\overline{B}_{st}=\frac{\tau_S}{1-e^{-\theta}}.
$$
在经典 CSMA 中，每次竞争后只发送一个数据包。
发送结束后节点立即释放信道并重新竞争。
因此其忙期 $B_{cl}$ 为常数，满足 $\overline{B}_{cl}=\tau_S$。
两种机制在竞争接入阶段采用相同的 LBT~与退避规则。
因此信道假期均值 $\overline{V}_c$ 可沿用前文结论：
\begin{equation}
	\overline{V}_c = \frac{p_{2}\tau_{F}+p_{0}\tau_{I}}{p_{1}},
\end{equation}
其中，$p_0 = e^{-\theta}$、$p_1 = \theta e^{-\theta}$ 和 $p_2 = 1 - p_0 - p_1$ 分别对应无传输、成功接入和发生碰撞的概率（$\theta = nq$）。

\section{吞吐量表达式与最优传输概率}

将 $\overline{B}_{cl}$ 与 $\overline{V}_c$ 代入吞吐量定义式，并引入有效传输修正因子 $\eta=(\tau_S-\tau_I)/\tau_S$，可得经典CSMA在饱和状态下的通用吞吐量表达式
\begin{equation}
	\label{eq:classic_csma_throughput}
	\hat{\lambda}_{\text{out}}^{\text{cl}} = \frac{\tau_{S}-\tau_{I}}{\tau_{S}} \cdot \frac{\overline{B}_{cl}}{\overline{B}_{cl}+\overline{V}_{c}} = \frac{\tau_{S}-\tau_{I}}{\tau_{S}+\frac{p_{2}\tau_{F}+p_{0}\tau_{I}}{p_{1}}}.
\end{equation}

为求理论容量上界，将 $\hat{\lambda}_{\text{out}}^{\text{cl}}$ 视为 $q$ 的函数。
求解最优化条件 $\frac{d\hat{\lambda}_{\text{out}}^{\text{cl}}}{dq}=0$ 后，得到最优传输概率
\begin{equation}
	\label{eq:optimal_q_classic}
	q^{*} = \frac{1 + W_{0}\left(\frac{\tau_{I}-\tau_{F}}{\tau_{F} e}\right)}{n},
\end{equation}
其中 $W_0(\cdot)$ 表示 Lambert W 函数主支\cite{2022_LambertW}。
将 $q^{*}$ 代回式 \eqref{eq:classic_csma_throughput}，即可得到正文图 \ref{fig:comparison_classic} 中经典CSMA曲线的理论上界。
